\documentclass{beamer}
\usepackage[slovene]{babel}
\usepackage[utf8]{inputenc}
\usepackage[T1]{fontenc}
\usepackage{lmodern}
\usepackage{amsmath, amssymb, bbm}

\usetheme{metropolis}
\beamertemplatenavigationsymbolsempty
\setbeamertemplate{caption}[numbered]

\usepackage{palatino}
\usefonttheme{serif}

% Krožne matrike
% Samo Primer
\institute[FMF]{Fakulteta za matematiko in fiziko}

% 
% Pozor: dokler ne dodate vsaj enega okolja za prosojnico, 
% se datoteka ne bo prevedla.
% 

\begin{document}




% začetek definicije
    Krožna matrika ?? je matrika oblike 

        % c_0  c_{n-1}    c_2  c_1 
        % c_1  c_0    c_{n-1}  c_2 
        % c_2  c_1    c_0  c_{n-1} 
        % 
        % c_{n-1}  c_{n-2}    c_1  c_0
        
% konec definicije



% začetek neoštevilčenega seznama
        % lastnost 1
        Lastni vektorji krožne matrike so 
        $v_j = (1, \omega_j, \omega_j^2, \ldots, \omega_j^{n-1})^T$, 
        kjer je ?? za $j=0, \ldots, n-1$.
        % lastnost 2
        Lastni vektorji krožne matrike so stolpci matrike enotske diskretne Fourierjeve transformacije, 
        ki jo lahko prikažemo kot Vandermondovo matriko.
        % lastnost 3
        Krožne matrike tvorijo komutativno algebro, ker je za poljubni dve matriki 
        $A$ in $B$, vsota $A + B$ tudi krožna matrika, prav tako je krožna matrika tudi njun produkt 
        $AB$, ter tudi velja $AB = BA$.
% konec neoštevilčenega seznama




\end{document}

